\documentclass{article}
\usepackage{mathtools} 
\usepackage{fontspec}
\usepackage[UTF8]{ctex}
\usepackage{amsthm}
\usepackage{mdframed}
\usepackage{xcolor}
\usepackage{amssymb}
\usepackage{amsmath}


% 定义新的带灰色背景的说明环境 zremark
\newmdtheoremenv[
  backgroundcolor=gray!10,
  % 边框与背景一致，边框线会消失
  linecolor=gray!10
]{zremark}{说明}

% 通用矩阵命令: \flexmatrix{矩阵名}{元素符号}{行数}{列数}
\newcommand{\flexmatrix}[4]{
  \[
  #1 = \begin{pmatrix}
    #2_{11}     & #2_{12}     & \cdots & #2_{1#4}   \\
    #2_{21}     & #2_{22}     & \cdots & #2_{2#4}   \\
    \vdots      & \vdots      & \ddots & \vdots     \\
    #2_{#31}    & #2_{#32}    & \cdots & #2_{#3#4}
  \end{pmatrix}
  \]
}

% 简化版命令（默认矩阵名为A，元素符号为a）: \quickmatrix{行数}{列数}
\newcommand{\quickmatrix}[2]{\flexmatrix{A}{a}{#1}{#2}}

\begin{document}
\title{注释 12.2}
\author{张志聪}
\maketitle

\begin{zremark}
  证明：命题12.7 推论2
\end{zremark}

(i)

对取定的正数$\epsilon = \frac{1}{2}|1 - q|$，由定理7.7可知，存在
$N > 0$，使得当$n > N$时，有
\begin{align*}
  \frac{u_{n + 1}}{u_n} < q + \epsilon
\end{align*}
当$q < 1$时，$q + \epsilon = \frac{1}{2} (1 + q) > 1$，于是
由定理12.7的(i)，推得级数$\sum u_n$是收敛的。

(ii)

同理可证。

\begin{zremark}
  证明：
  \begin{align*}
    \lim\limits_{n \to \infty} \sqrt[n]{1 + x^{2n}} = max\{1, x^2\}
  \end{align*}
  其中$x > 0$。
\end{zremark}

(i) $0 < x < 1$；

我们有
\begin{align*}
  x^{2n} < 1
\end{align*}
由对任意$a > 0$， $\lim\limits_{n \to \infty} \sqrt[n]{a} = 1$可知
\begin{align*}
  1 \leq \lim\limits_{n \to \infty} \sqrt[n]{1 + x^{2n}} \leq \lim\limits_{n \to \infty} \sqrt[n]{2} = 1
\end{align*}
由迫敛性可知
\begin{align*}
  \lim\limits_{n \to \infty} \sqrt[n]{1 + x^{2n}} = 1
\end{align*}

(ii) $x \geq 1$.

我们有
\begin{align*}
  x^{2n} \leq 1 + x^{2n} \leq 2x^{2n}
\end{align*}
所以
\begin{align*}
  \lim\limits_{n \to \infty} \sqrt[n]{x^{2n}} \leq \lim\limits_{n \to \infty} \sqrt[n]{1 + x^{2n}} \leq \lim\limits_{n \to \infty} \sqrt[n]{2x^{2n}} \\
  x^2 \leq \lim\limits_{n \to \infty} \sqrt[n]{1 + x^{2n}} \leq \lim\limits_{n \to \infty} \sqrt[n]{2} \cdot \lim\limits_{n \to \infty} \sqrt[n]{x^{2n}} = 1 \cdot x^2
\end{align*}

由迫敛性可知
\begin{align*}
  \lim\limits_{n \to \infty} \sqrt[n]{1 + x^{2n}} = x^2
\end{align*}

综上
\begin{equation*}
  \lim\limits_{n \to \infty} \sqrt[n]{1 + x^{2n}} =
  \begin{cases}
    1   & 0 < x < 1 \\
    x^2 & x \geq 1
  \end{cases}
\end{equation*}


\end{document}